\scanpage{239}
\mspage{113}
% archive sheet number: 238

\unsure{-|.} une \deleted{telle} factorisation de ce type pour $X\amalg X\xrightarrow{(\mathrm{id},\mathrm{id})}X$, i.e. un objet $I_X$ et une \deleted{\unsure{diagr.}} factorisation

\diagram{\begin{tikzcd}[ampersand replacement=\&,column sep=large]
X\amalg X\arrow[r,hook,"{i_X}","{\substack{\Vert\\(\delta_0^X,\delta_1^X)}}"']\arrow[rr,bend right=24,dashed,"{(\mathrm{id}_X,\mathrm{id}_X)}"'] \& I_X\arrow[r,"{p_X}"] \& X,
\end{tikzcd}}

avec $i_X\in C$, $p_X\in W$. On \deleted{\unsure{f|y}} associe alors à tout $f$ \deleted{\unsure{son|un}} \unsure{mapping cone|mapping-cone} $C(f)$ comme dans le diagramme (8.6.1) (où il faut lire $I_X$ ald $X\times I$), et on a alors $p_f\in W$, $i_f\in C$ comme voulu. Mais il faut supposer, pour savoir que $i_f\in C$, que \deleted{\unsure{NB|N.B.} \unsure{les diagrammes} \unsure{co--}} $X\in \unsure{\mathcal{M}|M}_C$ (objets $C$-cofibrants) i.e que pour \unsure{$\forall\,Y$|tt $Y$}, $X\unsure{\rightarrowtail|\longrightarrow}\deleted{\unsure{X}}Y\amalg X$ est dans $C$. (De plus, on a supposé tacitement $\mathcal{M}$ stable par $\unsure{\varinjlim|\underline{\lim}}$ finies.)

Le diagramme complet est

\diagram{\[
\left\{\begin{gathered}
\begin{tikzcd}[ampersand replacement=\&,column sep=large,row sep=large]
\& X\arrow[r,"f"]\arrow[d,"{\mathrm{inj}_0^X}"]\arrow[dr,phantom,"{\scriptstyle\mathrm{cocart}}" description] \& Y\arrow[d]\\
X\arrow[r,"{\mathrm{inj}_1^X}"]\arrow[dr,dashed,"{\delta_1^X}"'] \& X\amalg X\arrow[d,"{i_X}"]\arrow[r,"{f''}"]\arrow[dr,phantom,"{\scriptstyle\mathrm{cocart}}" description] \& Y\amalg X\arrow[d,"{\unsure{j|J}}"]\\
\& I_X\arrow[uu,bend right=20,dashed,"{p_X}"']\arrow[r,"{f'}"'] \& Z(f)\arrow[uu,bend right=38,dashed,"{p_f}"']
\end{tikzcd}\\
f=p_f i_f\ ,\ \text{où }\ i_f=f'\circ\delta_1^X=j\circ(\underbrace{f''\circ\mathrm{inj}_1^X}_{\lambda}) .
\end{gathered}\right.
\tag*{(8.\unsure{7}.1)}
\]}
